%%%%%%%%%%%%%%%%%%%%%%%%%%%%%%%%%%%%%%%%%
% Lachaise Assignment
% LaTeX Template
% Version 1.0 (26/6/2018)
%
% This template originates from:
% http://www.LaTeXTemplates.com
%
% Authors:
% Marion Lachaise & François Févotte
% Vel (vel@LaTeXTemplates.com)
%
% License:
% CC BY-NC-SA 3.0 (http://creativecommons.org/licenses/by-nc-sa/3.0/)
%
%%%%%%%%%%%%%%%%%%%%%%%%%%%%%%%%%%%%%%%%%

%----------------------------------------------------------------------------------------
%	PACKAGES AND OTHER DOCUMENT CONFIGURATIONS
%----------------------------------------------------------------------------------------

\documentclass{article}

\usepackage[margin=1cm, bottom=2cm]{geometry}

% \input{structure.tex} % Include the file specifying the document structure and custom commands

% \usepackage{indentfirst}\setlength{\parindent}{0em}
\usepackage{listings}
\usepackage{float}
\usepackage{amsmath,amsfonts,stmaryrd,amssymb} % Math packages
\usepackage{graphicx}
\usepackage{subcaption}
\usepackage{tikz}

% add linking w/ blue color
\usepackage{xcolor}
\usepackage{hyperref}

\hypersetup{
    colorlinks=true,
    linkcolor=blue,
    filecolor=magenta,
    urlcolor=blue,
}

\numberwithin{figure}{section}

\title{Draft}

\begin{document}
\maketitle

\section{Impact of frequency shift on amplitude}

We have derived that the numerical angular frequency is given by
\begin{equation}
    \omega_{\mathrm{ADI}}(k_x) = \frac{2}{\Delta t}\,\arctan\!\left(S\sin\frac{k_x\Delta x}{2}\right),
    \label{eq:adi-dispersion-relation}
\end{equation}
thus the frequency shift is given by
\begin{equation}
    \frac{\omega_{\mathrm{ADI}}(k_x)}{\omega_0} = \frac{2\arctan\!\left[S\sin(k\Delta z/2)\right]}{S\,k\Delta z}.
\end{equation}

Consider an input source whose temporal profile is given by
\begin{equation}
    s(t)= \alpha \exp\left[-\frac{(t-t_0)^2}{2T_P^2}\right] \sin(\omega_d t),
\end{equation}
compute the Fourier-transform by $\hat s(\omega)=\int_{-\infty}^{\infty}s(t)e^{-i\omega t}\,dt$,
\begin{equation}
    \hat s(\omega) = \frac{\alpha\sqrt{2\pi}T_P}{2i} e^{-i\omega t_0} \left[ e^{i\omega_d t_0} e^{-\frac12T_P^2(\omega-\omega_d)^2} - e^{-i\omega_d t_0} e^{-\frac12T_P^2(\omega+\omega_d)^2} \right].
\end{equation}
Consider $\omega_{\mathrm{ADI}}$ that is close to $\omega_d$, so we have
\begin{equation}
    e^{-\frac12T_P^2(\omega_{\mathrm{ADI}}+\omega_d)^2} \ll e^{-\frac12T_P^2(\omega_{\mathrm{ADI}}-\omega_d)^2},
\end{equation}
thus
\begin{equation}
    \hat s(\omega_{\mathrm{ADI}}) \approx \frac{\alpha\sqrt{2\pi}T_P}{2i}\, e^{-i(\omega_{\mathrm{ADI}}-\omega_d)t_0}\, e^{-\frac12 T_P^2(\omega_{\mathrm{ADI}}-\omega_d)^2}.
\end{equation}
Taking the magnitude,
\begin{equation}
    \bigl|\hat s(\omega_{\mathrm{ADI}})\bigr| \approx \frac{\alpha\sqrt{2\pi}T_P}{2}\, e^{-\frac12 T_P^2(\omega_{\mathrm{ADI}}-\omega_d)^2} \propto e^{-\frac12 T_P^2(\omega_{\mathrm{ADI}}-\omega_d)^2}.
\end{equation}

\section{Source term analysis}
\subsection{Electric field source terms}

Use the conventions $\delta_x^{E,2}=\delta_x^H\delta_x^E$ and $\delta_x^{H,2}=\delta_x^E\delta_x^H$.
Let $S_{E,1}^n$ be the electric field source term in between $n$ and $n+\frac{1}{2}$, and $S_{E,2}^n$ be that in between $n+\frac{1}{2}$ and $n+1$.
\begin{equation}
    \left(1-\gamma\delta_x^{E,2}\right)E_z^{n+\frac{1}{2}} = E_z^n+C_b\delta_x^H H_y^n+S_{E,1}^n \label{eq:implicit-electric-half-step}
\end{equation}
\begin{equation}
    H_y^{n+\frac{1}{2}} = H_y^n+D_b\delta_x^E E_z^{n+\frac{1}{2}} \label{eq:magnetic-half-step}
\end{equation}
\begin{equation}
    E_z^{n+1} = E_z^{n+\frac{1}{2}} + C_b\delta_x^H H_y^{n+\frac{1}{2}}+S_{E,2}^n \label{eq:electric-full-step}
\end{equation}
\begin{equation}
    H_y^{n+1} = H_y^{n+\frac{1}{2}} + D_b\delta_x^E E_z^{n+\frac{1}{2}} \label{eq:magnetic-full-step}
\end{equation}
Shifting back equations~\eqref{eq:electric-full-step} and \eqref{eq:magnetic-full-step} one full step gives
\begin{equation}
    E_z^n = E_z^{n-\frac{1}{2}} + C_b\delta_x^H H_y^{n-\frac{1}{2}}+S_{E,2}^{n-1} \label{eq:shifted-electric-step}
\end{equation}
\begin{equation}
    H_y^n = H_y^{n-\frac{1}{2}} + D_b\delta_x^E E_z^{n-\frac{1}{2}} \label{eq:shifted-magnetic-step}
\end{equation}
To eliminate $H_y$ from equations~\eqref{eq:implicit-electric-half-step}--\eqref{eq:electric-full-step}, substitute equation~\eqref{eq:magnetic-half-step} into equation~\eqref{eq:electric-full-step}:
\begin{equation}
    E_z^{n+1} = E_z^{n+\frac{1}{2}} + C_b\delta_x^H H_y^n + \gamma\delta_x^{E,2} E_z^{n+\frac{1}{2}}+S_{E,2}^n. \label{eq:e-after-h-substitution}
\end{equation}
Equation~\eqref{eq:implicit-electric-half-step} gives
\begin{equation}
    C_b\delta_x^H H_y^n = \left(1-\gamma\delta_x^{E,2}\right)E_z^{n+\frac{1}{2}}-E_z^n-S_{E,1}^n. \label{eq:h-elimination-identity}
\end{equation}
Using $\gamma\delta_x^{E,2}=C_bD_b\delta_x^H\delta_x^E$, equations~\eqref{eq:e-after-h-substitution} and \eqref{eq:h-elimination-identity} reduce to
\begin{equation}
    E_z^{n+\frac{1}{2}} = \frac{1}{2}\left(E_z^{n+1}+E_z^n+S_{E,1}^n-S_{E,2}^n\right). \label{eq:e-only-first-half}
\end{equation}
Shifting back one full step gives
\begin{equation}
    E_z^{n-\frac{1}{2}} = \frac{1}{2}\left(E_z^n+E_z^{n-1}+S_{E,1}^{n-1}-S_{E,2}^{n-1}\right). \label{eq:e-only-first-half-shifted}
\end{equation}
Next, eliminate $H_y$ using equations~\eqref{eq:implicit-electric-half-step}, \eqref{eq:shifted-electric-step}, and \eqref{eq:shifted-magnetic-step}. Applying $C_b\delta_x^H$ to both sides of equation~\eqref{eq:shifted-magnetic-step} gives
\begin{equation}
    C_b\delta_x^H H_y^n = C_b\delta_x^H H_y^{n-\frac{1}{2}}+\gamma\delta_x^{E,2}E_z^{n-\frac{1}{2}}. \label{eq:h-substitution-shifted}
\end{equation}
Eliminate $C_b\delta_x^H H_y^{n-\frac{1}{2}}$ by equation~\eqref{eq:shifted-electric-step} gives
\begin{equation}
    C_b\delta_x^H H_y^n=E_z^n-\left(1-\gamma\delta_x^{E,2}\right)E_z^{n-\frac{1}{2}}-S_{E,2}^{n-1}. \label{eq:h-elimination-identity-shifted}
\end{equation}
Substituting into equation~\eqref{eq:implicit-electric-half-step} to eliminate $H_y$:
\begin{equation}
    \left(1-\gamma\delta_x^{E,2}\right)\left(E_z^{n+\frac{1}{2}}+E_z^{n-\frac{1}{2}}\right) = 2E_z^n+S_{E,1}^n-S_{E,2}^{n-1}. \label{eq:e-only-centered}
\end{equation}
Finally, substitute equations~\eqref{eq:e-only-first-half} and \eqref{eq:e-only-first-half-shifted} into equation~\eqref{eq:e-only-centered} to eliminate the half-step values of $E_z$:
\begin{equation}
    \left(1-\gamma\delta_x^{E,2}\right)E_z^{n+1} = 2\left(1+\gamma\delta_x^{E,2}\right)E_z^n-\left(1-\gamma\delta_x^{E,2}\right)E_z^{n-1}+\left(1+\gamma\delta_x^{E,2}\right)S_{E,1}^n-\left(1-\gamma\delta_x^{E,2}\right)S_{E,1}^{n-1}+\left(1-\gamma\delta_x^{E,2}\right)S_{E,2}^n-\left(1+\gamma\delta_x^{E,2}\right)S_{E,2}^{n-1}. \label{eq:e-only-full-step}
\end{equation}
Recall that $\gamma=C_bD_b=c^2\Delta t^2/4$, then equation~\eqref{eq:e-only-full-step} results in,
\begin{align}
    \begin{split}
         & \frac{E_z^{n+1}-2E_z^n+E_z^{n-1}}{\Delta t^2}=                                                                                                                                                                                                                                                                                     \\
         & c^2\delta_x^{E,2}\left(\frac{E_z^{n+1}+2E_z^n+E_z^{n-1}}{4}\right)+\frac{1}{\Delta t^2}\left[\left(1+\gamma\delta_x^{E,2}\right)S_{E,1}^n-\left(1-\gamma\delta_x^{E,2}\right)S_{E,1}^{n-1}+\left(1-\gamma\delta_x^{E,2}\right)S_{E,2}^n-\left(1+\gamma\delta_x^{E,2}\right)S_{E,2}^{n-1}\right]. \label{eq:discrete-wave-equation}
    \end{split}
\end{align}

Suppose the same source is applied in both substeps, with $S_{E,1}^n=S_{E,2}^n=\frac{\Delta t}{2}S^{n+\frac{1}{2}}$ and $S_{E,1}^{n-1}=S_{E,2}^{n-1}=\frac{\Delta t}{2}S^{n-\frac{1}{2}}$. Then equation~\eqref{eq:discrete-wave-equation} becomes
\begin{equation}
    \frac{E_z^{n+1}-2E_z^n+E_z^{n-1}}{\Delta t^2}=c^2\delta_x^{E,2}\left(\frac{E_z^{n+1}+2E_z^n+E_z^{n-1}}{4}\right)+\frac{S^{n+\frac{1}{2}}-S^{n-\frac{1}{2}}}{\Delta t}. \label{eq:wave-equation-symmetric-source}
\end{equation}
Assume $S^{n+a}=S(t_n+a\Delta t)$ and $F^n=S_t(t_n)$. The source approximation in equation~\eqref{eq:wave-equation-symmetric-source} has the Taylor expansion
\begin{equation}
    \frac{S^{n+\frac{1}{2}}-S^{n-\frac{1}{2}}}{\Delta t}=F^n+\frac{\Delta t^2}{24}S_{ttt}^n+O(\Delta t^4), \qquad \mathrm{error}=\frac{\Delta t^2}{24}S_{ttt}^n+O(\Delta t^4). \label{eq:symmetric-source-error}
\end{equation}
Thus the two-substep midpoint-source case is second-order accurate for $F^n=S_t(t_n)$.

Alternatively, suppose the source is added only once in the first substep and evaluated at the midpoint, with $S_{E,1}^n=\Delta t S^{n+\frac{1}{2}}$, $S_{E,2}^n=0$, $S_{E,1}^{n-1}=\Delta t S^{n-\frac{1}{2}}$, and $S_{E,2}^{n-1}=0$. Then equation~\eqref{eq:discrete-wave-equation} becomes
\begin{equation}
    \frac{E_z^{n+1}-2E_z^n+E_z^{n-1}}{\Delta t^2}=c^2\delta_x^{E,2}\left(\frac{E_z^{n+1}+2E_z^n+E_z^{n-1}}{4}\right)+\frac{1}{\Delta t}\left[\left(1+\gamma\delta_x^{E,2}\right)S^{n+\frac{1}{2}}-\left(1-\gamma\delta_x^{E,2}\right)S^{n-\frac{1}{2}}\right]. \label{eq:wave-equation-single-midpoint-source}
\end{equation}
The source approximation in equation~\eqref{eq:wave-equation-single-midpoint-source} has the Taylor expansion
\begin{equation}
    F^n+\frac{2\gamma}{\Delta t}\delta_x^{E,2}S^n+\frac{\Delta t^2}{24}S_{ttt}^n+O(\Delta t^3). \label{eq:single-midpoint-source-error}
\end{equation}
Thus the single midpoint-source case is first-order accurate for a spatially varying $S$, and second-order when $\delta_x^{E,2}S=0$.

Alternatively, suppose the first source is evaluated at the first quarter-step and the second source at the third quarter-step:
$S_{E,1}^n=\frac{\Delta t}{2}S^{n+\frac{1}{4}}$, $S_{E,2}^n=\frac{\Delta t}{2}S^{n+\frac{3}{4}}$, $S_{E,1}^{n-1}=\frac{\Delta t}{2}S^{n-\frac{3}{4}}$, and $S_{E,2}^{n-1}=\frac{\Delta t}{2}S^{n-\frac{1}{4}}$. Then equation~\eqref{eq:discrete-wave-equation} becomes
\begin{equation}
    \frac{E_z^{n+1}-2E_z^n+E_z^{n-1}}{\Delta t^2}=c^2\delta_x^{E,2}\left(\frac{E_z^{n+1}+2E_z^n+E_z^{n-1}}{4}\right)+\frac{1}{2\Delta t}\left[\left(1+\gamma\delta_x^{E,2}\right)S^{n+\frac{1}{4}}-\left(1-\gamma\delta_x^{E,2}\right)S^{n-\frac{3}{4}}+\left(1-\gamma\delta_x^{E,2}\right)S^{n+\frac{3}{4}}-\left(1+\gamma\delta_x^{E,2}\right)S^{n-\frac{1}{4}}\right]. \label{eq:wave-equation-quarter-step-source}
\end{equation}
The source approximation in equation~\eqref{eq:wave-equation-quarter-step-source} has the Taylor expansion
\begin{equation}
    F^n-\frac{\gamma}{2}\delta_x^{E,2}F^n+\frac{7\Delta t^2}{96}S_{ttt}^n+O(\Delta t^4), \qquad \mathrm{error}=-\frac{\gamma}{2}\delta_x^{E,2}F^n+\frac{7\Delta t^2}{96}S_{ttt}^n+O(\Delta t^4). \label{eq:quarter-step-source-error}
\end{equation}
Since $\gamma=O(\Delta t^2)$, the quarter-step source case is also second-order accurate.

Assume $S(x,t)=\widehat{S}e^{i(\omega t-kx)}$, and define $\theta=\omega\Delta t$, $\delta_x^{E,2}e^{-ikx}=\lambda_xe^{-ikx}$, and $q=-\gamma\lambda_x\geq0$. Substituting into the source term on the right-hand side of equation~\eqref{eq:wave-equation-symmetric-source} gives
\begin{equation}
    \begin{aligned} F_{\mathrm{mid,2}}^n & =\frac{S^{n+\frac{1}{2}}-S^{n-\frac{1}{2}}}{\Delta t}                                      \\
                                     & =\frac{\widehat{S}e^{i(\omega t_n-kx)}}{\Delta t}\left(e^{i\theta/2}-e^{-i\theta/2}\right) \\
                                     & =\frac{2i\widehat{S}}{\Delta t}\sin\left(\frac{\theta}{2}\right)e^{i(\omega t_n-kx)},
    \end{aligned} \label{eq:midpoint-both-source-amplitude}
\end{equation}
so $\widehat{F}_{\mathrm{mid,2}}=\frac{2i\widehat{S}}{\Delta t}\sin\left(\frac{\theta}{2}\right)$.
Substituting into the source term on the right-hand side of equation~\eqref{eq:wave-equation-single-midpoint-source} gives
\begin{equation}
    \begin{aligned} F_{\mathrm{mid,1}}^n & =\frac{1}{\Delta t}\left[\left(1+\gamma\delta_x^{E,2}\right)S^{n+\frac{1}{2}}-\left(1-\gamma\delta_x^{E,2}\right)S^{n-\frac{1}{2}}\right] \\
                                     & =\frac{\widehat{S}e^{i(\omega t_n-kx)}}{\Delta t}\left[\left(1-q\right)e^{i\theta/2}-\left(1+q\right)e^{-i\theta/2}\right]                \\
                                     & =\frac{2\widehat{S}}{\Delta t}\left[i\sin\left(\frac{\theta}{2}\right)-q\cos\left(\frac{\theta}{2}\right)\right]e^{i(\omega t_n-kx)}.
    \end{aligned} \label{eq:midpoint-source-amplitude}
\end{equation}
so $\widehat{F}_{\mathrm{mid,1}}=\frac{2\widehat{S}}{\Delta t}\left[i\sin\left(\frac{\theta}{2}\right)-q\cos\left(\frac{\theta}{2}\right)\right]$.
Substituting into the source term on the right-hand side of equation~\eqref{eq:wave-equation-quarter-step-source} gives
\begin{equation}
    \begin{aligned} F_{\mathrm{quarter}}^n & =\frac{1}{2\Delta t}\left[\left(1+\gamma\delta_x^{E,2}\right)S^{n+\frac{1}{4}}-\left(1-\gamma\delta_x^{E,2}\right)S^{n-\frac{3}{4}}+\left(1-\gamma\delta_x^{E,2}\right)S^{n+\frac{3}{4}}-\left(1+\gamma\delta_x^{E,2}\right)S^{n-\frac{1}{4}}\right] \\
                                       & =\frac{\widehat{S}e^{i(\omega t_n-kx)}}{2\Delta t}\left[\left(1-q\right)e^{i\theta/4}-\left(1+q\right)e^{-3i\theta/4}+\left(1+q\right)e^{3i\theta/4}-\left(1-q\right)e^{-i\theta/4}\right]                                                           \\
                                       & =\frac{i\widehat{S}}{\Delta t}\left[\left(1-q\right)\sin\left(\frac{\theta}{4}\right)+\left(1+q\right)\sin\left(\frac{3\theta}{4}\right)\right]e^{i(\omega t_n-kx)}.
    \end{aligned} \label{eq:quarter-source-amplitude}
\end{equation}
so $\widehat{F}_{\mathrm{quarter}}=\frac{i\widehat{S}}{\Delta t}\left[\left(1-q\right)\sin\left(\frac{\theta}{4}\right)+\left(1+q\right)\sin\left(\frac{3\theta}{4}\right)\right]$.
Therefore the magnitude ratios among the midpoint and quarter-step source terms are
\begin{equation}
    \frac{\left|\widehat{F}_{\mathrm{mid,1}}\right|}{\left|\widehat{F}_{\mathrm{mid,2}}\right|}=\sqrt{1+q^{2}\cot^{2}\left(\frac{\theta}{2}\right)}, \label{eq:source-magnitude-ratio-mid}
\end{equation}
\begin{equation}
    \frac{\left|\widehat{F}_{\mathrm{quarter}}\right|}{\left|\widehat{F}_{\mathrm{mid,2}}\right|}=\left|\cos\left(\frac{\theta}{4}\right)+\frac{q\cos\left(\frac{\theta}{2}\right)}{2\cos\left(\frac{\theta}{4}\right)}\right|, \label{eq:source-magnitude-ratio-both}
\end{equation}
\begin{equation}
    \frac{\left|\widehat{F}_{\mathrm{quarter}}\right|}{\left|\widehat{F}_{\mathrm{mid,1}}\right|}=\frac{\left|\left(1-q\right)\sin\left(\frac{\theta}{4}\right)+\left(1+q\right)\sin\left(\frac{3\theta}{4}\right)\right|}{2\sqrt{\sin^{2}\left(\frac{\theta}{2}\right)+q^{2}\cos^{2}\left(\frac{\theta}{2}\right)}}. \label{eq:source-magnitude-ratio}
\end{equation}
For a spatially uniform source, $q=0$, both midpoint placements give the same magnitude, and the quarter-step magnitude is smaller by the factor
\begin{equation}
    \frac{\left|\widehat{F}_{\mathrm{mid,1}}\right|}{\left|\widehat{F}_{\mathrm{mid,2}}\right|}=1,\qquad
    \frac{\left|\widehat{F}_{\mathrm{quarter}}\right|}{\left|\widehat{F}_{\mathrm{mid,2}}\right|}=\frac{\left|\widehat{F}_{\mathrm{quarter}}\right|}{\left|\widehat{F}_{\mathrm{mid,1}}\right|}=\left|\cos\left(\frac{\omega\Delta t}{4}\right)\right|. \label{eq:uniform-source-magnitude-ratio}
\end{equation}

\subsection{Method of manufactured solution validation for electric source}

To test the electric-source placements of the previous section with a genuinely nonzero forcing, consider the one-dimensional wave equation
\begin{equation}
    E_{tt}-c^2 E_{xx}=f(x,t),\qquad x\in[0,L],
    \label{eq:mms-wave-equation}
\end{equation}
with PEC boundary condition, $E(0,t)=E(L,t)=0$. Choose a smooth manufactured solution,
\begin{equation}
    E_{\mathrm{MMS}}(x,t)=\sin(kx)\cos(\omega t),\qquad
    k=\frac{2\pi}{L},\qquad
    \omega=\frac{3\pi c}{L},
    \label{eq:mms-solution}
\end{equation}
where we set $\omega \neq c k$ in order to have a non-zero $f$.

Indeed,
\begin{equation}
    E_{tt}=-\omega^2\sin(kx)\cos(\omega t),\qquad
    E_{xx}=-k^2\sin(kx)\cos(\omega t),
\end{equation}
so
\begin{equation}
    f(x,t)=\bigl(c^2k^2-\omega^2\bigr)\sin(kx)\cos(\omega t)
    =-\frac{5\pi^2 c^2}{L^2}\sin(kx)\cos(\omega t).
    \label{eq:mms-forcing}
\end{equation}
In the ADI update, $\partial_t E=\varepsilon^{-1}\partial_x H+S$ with $F=S_t$ the wave-equation source of~\eqref{eq:discrete-wave-equation}. Matching $F=f$ therefore gives the continuous electric soft source
\begin{equation}
    S(x,t)=\frac{c^2k^2-\omega^2}{\omega}\sin(kx)\sin(\omega t).
    \label{eq:mms-soft-source}
\end{equation}
From~\eqref{eq:mms-solution}, we have:
\begin{equation}
    E(x,0)=\sin(kx),\qquad E_t(x,0)=0,
\end{equation}
and the compatible magnetic field is $H(x,0)=0$ (since $S(x,0)=0$ as well). The wavenumber $k=2\pi/L$ satisfies the PEC conditions.

To analyse how the magnitude will scale with respect to the CFL numbers, we model the system as a forced oscillation system with natural frequency $\omega_{\mathrm{ADI}}(k)$ and a forcing function whose frequency is $\omega_d=\omega$ from~\eqref{eq:mms-solution}.

We model the system as a simple undamped forced oscillator,
\begin{equation}
    a''+\omega_{\mathrm{ADI}}^2 a=F(t).
    \label{eq:forced-oscillator}
\end{equation}
First, we consider $F(t)=F_{\mathrm{mid,2}}$ from~\eqref{eq:midpoint-both-source-amplitude}.
With the manufactured soft source~\eqref{eq:mms-soft-source},
\begin{align}
    \begin{split}
        F_{\mathrm{mid,2}}
         & =\frac{S^{n+\frac{1}{2}}-S^{n-\frac{1}{2}}}{\Delta t}                                                                                                                                \\
         & =\frac{c^2k^2-\omega^2}{\omega}\sin(kx)\,
        \frac{\sin\bigl(\omega(t_n+\tfrac12\Delta t)\bigr)-\sin\bigl(\omega(t_n-\tfrac12\Delta t)\bigr)}{\Delta t}                                                                              \\
         & =\frac{c^2k^2-\omega^2}{\omega}\sin(kx)\,
        \frac{2}{\Delta t}\sin\left(\frac{\omega\Delta t}{2}\right)\cos(\omega t_n).
    \end{split}
    \label{eq:mms-F-mid2-real}
\end{align}
Thus corresponding system is
\begin{align}
    \begin{split}
        a''+\omega_{\mathrm{ADI}}^2 a
         & =\frac{2}{\Delta t}\sin\left(\frac{\omega\Delta t}{2}\right)\cdot\frac{c^2k^2-\omega^2}{\omega}\cos(\omega t), \\
        a(0)
         & =1,                                                                                                                                                                                  \\
        a'(0)
         & =0.
    \end{split}
    \label{eq:mms-forced-mid2}
\end{align}
Solving this system yields
\begin{equation}
    \begin{aligned}
        a(t)
         & =\left(1-A_p\right)\cos(\omega_{\mathrm{ADI}}t)+A_p\cos(\omega t), \\
        A_p
         & =\mathrm{sinc}\left(\frac{\omega\Delta t}{2}\right)\frac{c^2k^2-\omega^2}{\omega_{\mathrm{ADI}}^2-\omega^2}.
    \end{aligned}
    \label{eq:mms-solution-mid2}
\end{equation}

Next, consider $F(t)=F_{\mathrm{mid,1}}$ from~\eqref{eq:midpoint-source-amplitude}.
On the mode $\sin(kx)$ one has $\delta_x^{E,2}\sin(kx)=\lambda_x\sin(kx)$ and $q=-\gamma\lambda_x$, so with~\eqref{eq:mms-soft-source},
\begin{align}
    \begin{split}
        F_{\mathrm{mid,1}}
         & =\frac{1}{\Delta t}\left[\left(1+\gamma\delta_x^{E,2}\right)S^{n+\frac{1}{2}}-\left(1-\gamma\delta_x^{E,2}\right)S^{n-\frac{1}{2}}\right]                                                                                       \\
         & =\frac{c^2k^2-\omega^2}{\omega\Delta t}\sin(kx)\Bigl[
            (1-q)\sin\bigl(\omega(t_n+\tfrac12\Delta t)\bigr)
            -(1+q)\sin\bigl(\omega(t_n-\tfrac12\Delta t)\bigr)
            \Bigr]                                                                                                                                                                                                              \\
         & =\frac{2(c^2k^2-\omega^2)}{\omega\Delta t}\sin(kx)\left[
            \sin\left(\frac{\omega\Delta t}{2}\right)\cos(\omega t_n)
            -q\cos\left(\frac{\omega\Delta t}{2}\right)\sin(\omega t_n)
            \right].
    \end{split}
    \label{eq:mms-F-mid1-real}
\end{align}
The corresponding modal system is
\begin{align}
    \begin{split}
        a''+\omega_{\mathrm{ADI}}^2 a
         & =\frac{2(c^2k^2-\omega^2)}{\omega\Delta t}\left[
            \sin\left(\frac{\omega\Delta t}{2}\right)\cos(\omega t)
            -q\cos\left(\frac{\omega\Delta t}{2}\right)\sin(\omega t)
            \right], \\
        a(0)
         & =1,                                                                                                                                                                                                                                                                                   \\
        a'(0)
         & =0.
    \end{split}
    \label{eq:mms-forced-mid1}
\end{align}
Writing the drive as $A\cos(\omega t)+B\sin(\omega t)$ with
\begin{equation}
    A=\frac{2(c^2k^2-\omega^2)}{\omega\Delta t}\sin\left(\frac{\omega\Delta t}{2}\right),\qquad
    B=-\frac{2(c^2k^2-\omega^2)}{\omega\Delta t}\,q\cos\left(\frac{\omega\Delta t}{2}\right),
    \label{eq:mms-mid1-AB}
\end{equation}
the solution is
\begin{equation}
    \begin{aligned}
        a(t)
         & =\left(1-\frac{A}{\omega_{\mathrm{ADI}}^2-\omega^2}\right)\cos(\omega_{\mathrm{ADI}}t)
        -\frac{\omega}{\omega_{\mathrm{ADI}}}\frac{B}{\omega_{\mathrm{ADI}}^2-\omega^2}\sin(\omega_{\mathrm{ADI}}t) \\
         & \quad+\frac{A}{\omega_{\mathrm{ADI}}^2-\omega^2}\cos(\omega t)
        +\frac{B}{\omega_{\mathrm{ADI}}^2-\omega^2}\sin(\omega t).
    \end{aligned}
    \label{eq:mms-solution-mid1}
\end{equation}

Finally, consider $F(t)=F_{\mathrm{quarter}}$ from~\eqref{eq:quarter-source-amplitude}.
With~\eqref{eq:mms-soft-source},
\begin{align}
    \begin{split}
        F_{\mathrm{quarter}}
         & =\frac{1}{2\Delta t}\Bigl[
            \left(1+\gamma\delta_x^{E,2}\right)S^{n+\frac{1}{4}}
            -\left(1-\gamma\delta_x^{E,2}\right)S^{n-\frac{3}{4}}                                                                                                                                                                \\
         & \qquad
            +\left(1-\gamma\delta_x^{E,2}\right)S^{n+\frac{3}{4}}
            -\left(1+\gamma\delta_x^{E,2}\right)S^{n-\frac{1}{4}}
            \Bigr]                                                                                                                                                                                                              \\
         & =\frac{c^2k^2-\omega^2}{\omega\Delta t}\sin(kx)\left[
            (1-q)\sin\left(\frac{\omega\Delta t}{4}\right)
            +(1+q)\sin\left(\frac{3\omega\Delta t}{4}\right)
            \right]\cos(\omega t_n).
    \end{split}
    \label{eq:mms-F-quarter-real}
\end{align}
The corresponding modal system is
\begin{align}
    \begin{split}
        a''+\omega_{\mathrm{ADI}}^2 a
         & =\frac{c^2k^2-\omega^2}{\omega\Delta t}\left[
            (1-q)\sin\left(\frac{\omega\Delta t}{4}\right)
            +(1+q)\sin\left(\frac{3\omega\Delta t}{4}\right)
            \right]\cos(\omega t), \\
        a(0)
         & =1,                                                                                                                                                                                                                                                                                   \\
        a'(0)
         & =0.
    \end{split}
    \label{eq:mms-forced-quarter}
\end{align}
Solving yields
\begin{equation}
    \begin{aligned}
        a(t)
         & =\left(1-A_p\right)\cos(\omega_{\mathrm{ADI}}t)+A_p\cos(\omega t), \\
        A_p
         & =\frac{c^2k^2-\omega^2}{(\omega_{\mathrm{ADI}}^2-\omega^2)\,\omega\Delta t}\left[
            (1-q)\sin\left(\frac{\omega\Delta t}{4}\right)
            +(1+q)\sin\left(\frac{3\omega\Delta t}{4}\right)
            \right].
    \end{aligned}
    \label{eq:mms-solution-quarter}
\end{equation}
Relative to mid,2, the particular amplitudes satisfy the same magnitude ratios as~\eqref{eq:source-magnitude-ratio-mid}--\eqref{eq:source-magnitude-ratio-both}.


\subsection{Magnetic field source terms}
Let $S_{H,1}^n$ be the magnetic field source term in between $n$ and $n+\frac{1}{2}$, and $S_{H,2}^n$ be that in between $n+\frac{1}{2}$ and $n+1$.
\begin{equation}
    \left(1-\gamma\delta_x^{E,2}\right)E_z^{n+\frac{1}{2}}=E_z^n+C_b\delta_x^H H_y^n. \label{eq:hsource-implicit-electric-half-step}
\end{equation}
\begin{equation}
    H_y^{n+\frac{1}{2}}=H_y^n+D_b\delta_x^E E_z^{n+\frac{1}{2}}+S_{H,1}^n. \label{eq:hsource-magnetic-half-step}
\end{equation}
\begin{equation}
    E_z^{n+1}=E_z^{n+\frac{1}{2}}+C_b\delta_x^H H_y^{n+\frac{1}{2}}. \label{eq:hsource-electric-full-step}
\end{equation}
\begin{equation}
    H_y^{n+1}=H_y^{n+\frac{1}{2}}+D_b\delta_x^E E_z^{n+\frac{1}{2}}+S_{H,2}^n. \label{eq:hsource-magnetic-full-step}
\end{equation}
Shifting back equations~\eqref{eq:hsource-electric-full-step} and \eqref{eq:hsource-magnetic-full-step} one full step gives
\begin{equation}
    E_z^n=E_z^{n-\frac{1}{2}}+C_b\delta_x^H H_y^{n-\frac{1}{2}}. \label{eq:hsource-shifted-electric-step}
\end{equation}
\begin{equation}
    H_y^n=H_y^{n-\frac{1}{2}}+D_b\delta_x^E E_z^{n-\frac{1}{2}}+S_{H,2}^{n-1}. \label{eq:hsource-shifted-magnetic-step}
\end{equation}
To eliminate $E_z$ from equations~\eqref{eq:hsource-magnetic-half-step}--\eqref{eq:hsource-magnetic-full-step}, equation~\eqref{eq:hsource-magnetic-half-step} gives
\begin{equation}
    D_b\delta_x^E E_z^{n+\frac{1}{2}}=H_y^{n+\frac{1}{2}}-H_y^n-S_{H,1}^n. \label{eq:hsource-e-elimination-identity}
\end{equation}
Substituting into equation~\eqref{eq:hsource-magnetic-full-step} gives
\begin{equation}
    H_y^{n+1}=2H_y^{n+\frac{1}{2}}-H_y^n-S_{H,1}^n+S_{H,2}^n. \label{eq:hsource-h-after-e-substitution}
\end{equation}
Solving for the half-step value gives
\begin{equation}
    H_y^{n+\frac{1}{2}}=\frac{1}{2}\left(H_y^{n+1}+H_y^n+S_{H,1}^n-S_{H,2}^n\right). \label{eq:hsource-h-only-first-half}
\end{equation}
Shifting back one full step gives
\begin{equation}
    H_y^{n-\frac{1}{2}}=\frac{1}{2}\left(H_y^n+H_y^{n-1}+S_{H,1}^{n-1}-S_{H,2}^{n-1}\right). \label{eq:hsource-h-only-first-half-shifted}
\end{equation}
Next, eliminate $E_z$ using equations~\eqref{eq:hsource-implicit-electric-half-step}, \eqref{eq:hsource-magnetic-half-step}, \eqref{eq:hsource-shifted-electric-step}, and \eqref{eq:hsource-shifted-magnetic-step}. Applying $D_b\delta_x^E$ to both sides of equation~\eqref{eq:hsource-shifted-electric-step} gives
\begin{equation}
    D_b\delta_x^E E_z^n=D_b\delta_x^E E_z^{n-\frac{1}{2}}+\gamma\delta_x^{H,2}H_y^{n-\frac{1}{2}}. \label{eq:hsource-e-substitution-shifted}
\end{equation}
Eliminate $D_b\delta_x^E E_z^{n-\frac{1}{2}}$ by equation~\eqref{eq:hsource-shifted-magnetic-step} gives
\begin{equation}
    D_b\delta_x^E E_z^n=H_y^n-\left(1-\gamma\delta_x^{H,2}\right)H_y^{n-\frac{1}{2}}-S_{H,2}^{n-1}. \label{eq:hsource-e-elimination-identity-shifted}
\end{equation}




Applying $D_b\delta_x^E$ to both sides of equation~\eqref{eq:hsource-implicit-electric-half-step} gives
\begin{equation}
    \left(1-\gamma\delta_x^{H,2}\right)D_b\delta_x^E E_z^{n+\frac{1}{2}}=D_b\delta_x^E E_z^n+\gamma\delta_x^{H,2}H_y^n. \label{eq:hsource-applied-implicit-electric-step}
\end{equation}
Eliminate $D_b\delta_x^E E_z^{n+\frac{1}{2}}$ by equation~\eqref{eq:hsource-magnetic-half-step} gives
\begin{equation}
    D_b\delta_x^E E_z^n=\left(1-\gamma\delta_x^{H,2}\right)H_y^{n+\frac{1}{2}}-H_y^n-\left(1-\gamma\delta_x^{H,2}\right)S_{H,1}^n. \label{eq:hsource-e-n-elimination-identity}
\end{equation}
Equating equations~\eqref{eq:hsource-e-n-elimination-identity} and \eqref{eq:hsource-e-elimination-identity-shifted} eliminates $E_z^n$:
\begin{equation}
    \left(1-\gamma\delta_x^{H,2}\right)\left(H_y^{n+\frac{1}{2}}+H_y^{n-\frac{1}{2}}\right)=2H_y^n+\left(1-\gamma\delta_x^{H,2}\right)S_{H,1}^n-S_{H,2}^{n-1}. \label{eq:hsource-h-only-centered}
\end{equation}


Finally, substitute equations~\eqref{eq:hsource-h-only-first-half} and \eqref{eq:hsource-h-only-first-half-shifted} into equation~\eqref{eq:hsource-h-only-centered} to eliminate the half-step values of $H_y$:
\begin{equation}
    \left(1-\gamma\delta_x^{H,2}\right)H_y^{n+1}=2\left(1+\gamma\delta_x^{H,2}\right)H_y^n-\left(1-\gamma\delta_x^{H,2}\right)H_y^{n-1}+\left(1-\gamma\delta_x^{H,2}\right)S_{H,1}^n-\left(1-\gamma\delta_x^{H,2}\right)S_{H,1}^{n-1}+\left(1-\gamma\delta_x^{H,2}\right)S_{H,2}^n-\left(1+\gamma\delta_x^{H,2}\right)S_{H,2}^{n-1}. \label{eq:hsource-h-only-full-step}
\end{equation}
Recall that $\gamma=C_bD_b=c^2\Delta t^2/4$, then equation~\eqref{eq:hsource-h-only-full-step} results in,
\begin{align}
    \begin{split}
         & \frac{H_y^{n+1}-2H_y^n+H_y^{n-1}}{\Delta t^2}=                                                                                                                                                                                                                                                                                             \\
         & c^2\delta_x^{H,2}\left(\frac{H_y^{n+1}+2H_y^n+H_y^{n-1}}{4}\right)+\frac{1}{\Delta t^2}\left[\left(1-\gamma\delta_x^{H,2}\right)S_{H,1}^n-\left(1-\gamma\delta_x^{H,2}\right)S_{H,1}^{n-1}+\left(1-\gamma\delta_x^{H,2}\right)S_{H,2}^n-\left(1+\gamma\delta_x^{H,2}\right)S_{H,2}^{n-1}\right]. \label{eq:hsource-discrete-wave-equation}
    \end{split}
\end{align}

Suppose the same source is applied in both substeps, with $S_{H,1}^n=S_{H,2}^n=\frac{\Delta t}{2}S_H^{n+\frac{1}{2}}$ and $S_{H,1}^{n-1}=S_{H,2}^{n-1}=\frac{\Delta t}{2}S_H^{n-\frac{1}{2}}$. Then equation~\eqref{eq:hsource-discrete-wave-equation} becomes
\begin{equation}
    \frac{H_y^{n+1}-2H_y^n+H_y^{n-1}}{\Delta t^2}=c^2\delta_x^{H,2}\left(\frac{H_y^{n+1}+2H_y^n+H_y^{n-1}}{4}\right)+\frac{\left(1-\gamma\delta_x^{H,2}\right)S_H^{n+\frac{1}{2}}-S_H^{n-\frac{1}{2}}}{\Delta t}. \label{eq:hsource-wave-equation-symmetric-source}
\end{equation}
Assume $S_H^{n+a}=S_H(t_n+a\Delta t)$ and $F_H^n=(S_H)_t(t_n)$. The source approximation in equation~\eqref{eq:hsource-wave-equation-symmetric-source} has the Taylor expansion
\begin{equation}
    F_H^n-\frac{\gamma}{\Delta t}\delta_x^{H,2}S_H^n-\frac{\gamma}{2}\delta_x^{H,2}F_H^n+\frac{\Delta t^2}{24}(S_H)_{ttt}^n+O(\Delta t^3). \label{eq:hsource-symmetric-source-error}
\end{equation}
Thus the two-substep midpoint-source case is first-order accurate for a spatially varying $S_H$, and second-order when $\delta_x^{H,2}S_H=0$.

Alternatively, suppose the source is added only once in the first substep and evaluated at the midpoint, with $S_{H,1}^n=\Delta t S_H^{n+\frac{1}{2}}$, $S_{H,2}^n=0$, $S_{H,1}^{n-1}=\Delta t S_H^{n-\frac{1}{2}}$, and $S_{H,2}^{n-1}=0$. Then the source terms in equation~\eqref{eq:hsource-discrete-wave-equation} reduce to
\begin{equation}
    \Delta t\left(1-\gamma\delta_x^{H,2}\right)\left(S_H^{n+\frac{1}{2}}-S_H^{n-\frac{1}{2}}\right). \label{eq:hsource-single-midpoint-source-reduction}
\end{equation}
Therefore equation~\eqref{eq:hsource-discrete-wave-equation} becomes
\begin{equation}
    \frac{H_y^{n+1}-2H_y^n+H_y^{n-1}}{\Delta t^2}=c^2\delta_x^{H,2}\left(\frac{H_y^{n+1}+2H_y^n+H_y^{n-1}}{4}\right)+\left(1-\gamma\delta_x^{H,2}\right)\frac{S_H^{n+\frac{1}{2}}-S_H^{n-\frac{1}{2}}}{\Delta t}. \label{eq:hsource-wave-equation-single-midpoint-source}
\end{equation}

The source approximation in equation~\eqref{eq:hsource-wave-equation-single-midpoint-source} has the Taylor expansion
\begin{equation}
    F_H^n-\gamma\delta_x^{H,2}F_H^n+\frac{\Delta t^2}{24}(S_H)_{ttt}^n+O(\Delta t^4). \label{eq:hsource-single-midpoint-source-error}
\end{equation}
Since $\gamma=O(\Delta t^2)$, the single midpoint-source case is second-order accurate for $F_H^n=(S_H)_t(t_n)$.

Alternatively, suppose the first source is evaluated at the first quarter-step and the second source at the third quarter-step:
$S_{H,1}^n=\frac{\Delta t}{2}S_H^{n+\frac{1}{4}}$, $S_{H,2}^n=\frac{\Delta t}{2}S_H^{n+\frac{3}{4}}$, $S_{H,1}^{n-1}=\frac{\Delta t}{2}S_H^{n-\frac{3}{4}}$, and $S_{H,2}^{n-1}=\frac{\Delta t}{2}S_H^{n-\frac{1}{4}}$. Then equation~\eqref{eq:hsource-discrete-wave-equation} becomes
\begin{align}
    \begin{split}
         & \frac{H_y^{n+1}-2H_y^n+H_y^{n-1}}{\Delta t^2}=                                                                                                                                                                                                                                                                                                                                       \\
         & c^2\delta_x^{H,2}\left(\frac{H_y^{n+1}+2H_y^n+H_y^{n-1}}{4}\right)+\frac{1}{2\Delta t}\left[\left(1-\gamma\delta_x^{H,2}\right)S_H^{n+\frac{1}{4}}-\left(1-\gamma\delta_x^{H,2}\right)S_H^{n-\frac{3}{4}}+\left(1-\gamma\delta_x^{H,2}\right)S_H^{n+\frac{3}{4}}-\left(1+\gamma\delta_x^{H,2}\right)S_H^{n-\frac{1}{4}}\right]. \label{eq:hsource-wave-equation-quarter-step-source}
    \end{split}
\end{align}
The source approximation in equation~\eqref{eq:hsource-wave-equation-quarter-step-source} has the Taylor expansion
\begin{equation}
    F_H^n-\frac{\gamma}{\Delta t}\delta_x^{H,2}S_H^n-\frac{3\gamma}{4}\delta_x^{H,2}F_H^n+\frac{7\Delta t^2}{96}(S_H)_{ttt}^n+O(\Delta t^3). \label{eq:hsource-quarter-step-source-error}
\end{equation}
Thus the quarter-step source case is however first-order accurate for a spatially varying $S_H$, and second-order when $\delta_x^{H,2}S_H=0$.

Assume $S_H(x,t)=\widehat{S}_H e^{i(\omega t-kx)}$, and define $\theta=\omega\Delta t$, $\delta_x^{H,2}e^{-ikx}=\lambda_{H,x}e^{-ikx}$, and $q_H=-\gamma\lambda_{H,x}\geq0$. Substituting into the source term on the right-hand side of equation~\eqref{eq:hsource-wave-equation-symmetric-source} gives
\begin{equation}
    \begin{aligned} F_{H,\mathrm{mid,2}}^n & =\frac{\left(1-\gamma\delta_x^{H,2}\right)S_H^{n+\frac{1}{2}}-S_H^{n-\frac{1}{2}}}{\Delta t}                                                               \\
                                       & =\frac{\widehat{S}_H e^{i(\omega t_n-kx)}}{\Delta t}\left[\left(1+q_H\right)e^{i\theta/2}-e^{-i\theta/2}\right]                                            \\
                                       & =\frac{\widehat{S}_H}{\Delta t}\left[q_H\cos\left(\frac{\theta}{2}\right)+i\left(2+q_H\right)\sin\left(\frac{\theta}{2}\right)\right]e^{i(\omega t_n-kx)}.
    \end{aligned} \label{eq:hsource-midpoint-both-source-amplitude}
\end{equation}
so $\widehat{F}_{H,\mathrm{mid,2}}=\frac{\widehat{S}_H}{\Delta t}\left[q_H\cos\left(\frac{\theta}{2}\right)+i\left(2+q_H\right)\sin\left(\frac{\theta}{2}\right)\right]$.
Substituting into the source term on the right-hand side of equation~\eqref{eq:hsource-wave-equation-single-midpoint-source} gives
\begin{equation}
    \begin{aligned} F_{H,\mathrm{mid,1}}^n & =\left(1-\gamma\delta_x^{H,2}\right)\frac{S_H^{n+\frac{1}{2}}-S_H^{n-\frac{1}{2}}}{\Delta t}                    \\
                                       & =\frac{\left(1+q_H\right)\widehat{S}_H e^{i(\omega t_n-kx)}}{\Delta t}\left(e^{i\theta/2}-e^{-i\theta/2}\right) \\
                                       & =\frac{2i\left(1+q_H\right)\widehat{S}_H}{\Delta t}\sin\left(\frac{\theta}{2}\right)e^{i(\omega t_n-kx)}.
    \end{aligned} \label{eq:hsource-midpoint-source-amplitude}
\end{equation}
so $\widehat{F}_{H,\mathrm{mid,1}}=\frac{2i\left(1+q_H\right)\widehat{S}_H}{\Delta t}\sin\left(\frac{\theta}{2}\right)$.
Substituting into the source term on the right-hand side of equation~\eqref{eq:hsource-wave-equation-quarter-step-source} gives
\begin{equation}
    \begin{aligned} F_{H,\mathrm{quarter}}^n & =\frac{1}{2\Delta t}\left[\left(1-\gamma\delta_x^{H,2}\right)S_H^{n+\frac{1}{4}}-\left(1-\gamma\delta_x^{H,2}\right)S_H^{n-\frac{3}{4}}+\left(1-\gamma\delta_x^{H,2}\right)S_H^{n+\frac{3}{4}}-\left(1+\gamma\delta_x^{H,2}\right)S_H^{n-\frac{1}{4}}\right] \\
                                         & =\frac{\widehat{S}_H e^{i(\omega t_n-kx)}}{2\Delta t}\left[\left(1+q_H\right)e^{i\theta/4}-\left(1+q_H\right)e^{-3i\theta/4}+\left(1+q_H\right)e^{3i\theta/4}-\left(1-q_H\right)e^{-i\theta/4}\right]                                                        \\
                                         & =\frac{\widehat{S}_H}{\Delta t}\left[q_H\cos\left(\frac{\theta}{4}\right)+i\left(\sin\left(\frac{\theta}{4}\right)+\left(1+q_H\right)\sin\left(\frac{3\theta}{4}\right)\right)\right]e^{i(\omega t_n-kx)}.
    \end{aligned} \label{eq:hsource-quarter-source-amplitude}
\end{equation}
Therefore the magnitude ratios between the quarter-step and midpoint magnetic-source terms are
\begin{equation}
    \frac{\left|\widehat{F}_{H,\mathrm{quarter}}\right|}{\left|\widehat{F}_{H,\mathrm{mid,2}}\right|}=\frac{\sqrt{q_H^{2}\cos^{2}\left(\frac{\theta}{4}\right)+\left[\sin\left(\frac{\theta}{4}\right)+\left(1+q_H\right)\sin\left(\frac{3\theta}{4}\right)\right]^{2}}}{\sqrt{q_H^{2}\cos^{2}\left(\frac{\theta}{2}\right)+\left(2+q_H\right)^{2}\sin^{2}\left(\frac{\theta}{2}\right)}}, \label{eq:hsource-source-magnitude-ratio-both}
\end{equation}
\begin{equation}
    \frac{\left|\widehat{F}_{H,\mathrm{quarter}}\right|}{\left|\widehat{F}_{H,\mathrm{mid,1}}\right|}=\frac{\sqrt{q_H^2\cos^2\left(\frac{\theta}{4}\right)+\left[\sin\left(\frac{\theta}{4}\right)+\left(1+q_H\right)\sin\left(\frac{3\theta}{4}\right)\right]^2}}{2\left(1+q_H\right)\left|\sin\left(\frac{\theta}{2}\right)\right|}. \label{eq:hsource-source-magnitude-ratio}
\end{equation}
For a spatially uniform source, $q_H=0$, both midpoint placements give the same magnitude, and the quarter-step magnitude is smaller by the factor
\begin{equation}
    \frac{\left|\widehat{F}_{H,\mathrm{quarter}}\right|}{\left|\widehat{F}_{H,\mathrm{mid,2}}\right|}=\frac{\left|\widehat{F}_{H,\mathrm{quarter}}\right|}{\left|\widehat{F}_{H,\mathrm{mid,1}}\right|}=\left|\cos\left(\frac{\omega\Delta t}{4}\right)\right|. \label{eq:hsource-uniform-source-magnitude-ratio}
\end{equation}

\end{document}